\documentclass[11pt,a4paper]{article}

% --- Kodowanie i język ---
\usepackage[utf8]{inputenc}
\usepackage[T1]{fontenc}
\usepackage[polish]{babel}
\usepackage{lmodern}

% --- Układ strony ---
\usepackage[a4paper,margin=2.2cm]{geometry}
\usepackage{microtype}
\usepackage{parskip}

% --- Listy, tabele, linki ---
\usepackage{enumitem}
\usepackage{longtable}
\usepackage{array}
\usepackage{booktabs}
\usepackage[hidelinks]{hyperref}
\usepackage{xurl}

% --- Bibliografia (biblatex + biber, plik references.bib) ---
\usepackage{csquotes}
\usepackage[backend=biber,style=numeric-comp,sorting=none,
            maxbibnames=99,giveninits=true,isbn=false]{biblatex}
\addbibresource{references.bib}
% Nie drukuj dat dostępu z Zotero ("Term. wiz.")
\AtEveryBibitem{\clearfield{urlyear}\clearfield{urlmonth}\clearfield{urlday}}
% Usuń pole volume, gdy jest tylko rokiem (np. ICLR z Zotero: "T. 2025")
\DeclareSourcemap{\maps[datatype=bibtex]{
  \map{\step[fieldsource=volume, match=\regexp{\A(19|20)\d\d\Z}, final]
        \step[fieldset=volume, null]}
  % Nie drukuj pola language ("en.") eksportowanego przez Zotero
  \map{\step[fieldset=language, null]}
  % Jeśli praca ma DOI, nie drukuj dodatkowo URL
  \map{\step[fieldsource=doi, final]
        \step[fieldset=url, null]}}}

\newcolumntype{L}[1]{>{\raggedright\arraybackslash}p{#1}}

\title{Etap I -- Studium literaturowe}
\author{%
  \small Miłosz Bedryło\\\small 272934 \and
  \small Grzegorz Sienkiewicz\\\small 272868 \and
  \small Kamil Sromek\\\small 272903
}
\date{\small Grupa 3}

\begin{document}

\maketitle
\vspace{-1.5em}

\section*{Obszar i problem}
Wraz z dynamicznym rozwojem technologii w zakresie sztucznej inteligencji, w tym dużych modeli językowych (ang. \emph{Large Language Models - LLMs}), rozwiązania oparte na nich coraz częściej znajdują zastosowanie również w nauce \cite{guetal2026a}. Przy wyborze modelu wykorzystywanego do pracy koszt nie jest jedynym czynnikiem wpływającym na wybór modelu, innym istotnym czynnikiem jest ocena modelu. W naszym projekcie położymy nacisk na niektóre z trudności związane właśnie z procesem oceny takich modeli.

Przykładowym sposobem oceny modeli LLM jest wykorzystanie zbioru złożonego z pytań oraz odpowiadających im wzorcowych odpowiedzi. Mając taki zestaw, model ocenia się, wysyłając do niego pytania ze zbioru, a następnie porównując uzyskane odpowiedzi z wzorcowymi \cite{zhengetal2023c,chenetal2024a}. Podejście to jest konceptualnie proste, jednak chcąc rzetelnie przeprowadzić badanie jednym z problemów staje się ilość czasu potrzebna na samo odpytywanie modelu, wynika to nie tylko z wielkości zbioru testowego, ale też z potrzeby powtarzania zapytań w celu uzyskania istotnych statystycznie wyników \cite{yagubyan2026,lee2026how}. 

Jednym z zaproponowanych rozwiązań powyższego problemu jest zastąpienie ludzkiego badacza w procesie oceny modelu modelem LLM spełniającym to samo zadanie. Takie właśnie podejście w literaturze opisywane jest jako \emph{LLM-as-a-judge} \cite{zhengetal2023c,guetal2026a}. Metoda ta w wielu przypadkach pozwala nie tylko na oszczędność czasu, ale co jest z tym powiązane, również i kosztów związanych z procesem (przykładowo, wynagrodzeniem dla osób przeprowadzających badanie) \cite{zhengetal2023c,wangetal2024}.

Wykorzystanie modeli LLM jako sędziów oceniających inne modele jest też obarczone pewnym ryzykiem. Modele te, podobnie jak ludzie, również narażone są na popełnianie błędów poznawczych \cite{kooetal2024a,ICLR2025_fdca08d3,chenetal2024a}, za przykłady można podać faworyzowanie pierwszej odpowiedzi \cite{wangetal2024,shietal2025a}, przychylniejszą ocenę dłuższych tekstów \cite{zhengetal2023c,parketal2024a,soumik2026} czy też bias powodujący zawyżanie ocen w przypadku kiedy model jest oceniającym i ocenianym zarazem \cite{panicksseryetal2024b,wataokaetal2025,yangetal2026a}.

Często spotykanym w literaturze modelem jest GPT-4 \cite{zhengetal2023c,wangetal2024,ICLR2024_afc8b034,panicksseryetal2024b}, natomiast będąc zamkniętym modelem, wykorzystywanie go wiąże się z dodatkowymi kosztami. Jedną z alternatyw dla GPT są modele otwarte, które również można wykorzystać do tego zadania \cite{kooetal2024a,normanetal2026,yangetal2026,zhouetal2026}.

W badaniach opisujących wykorzystanie otwartych modeli jako sędziów, nie badano wpływu kwantyzacji, a jedynie w \cite{angheletal2026} wszyscy sędziowie byli uruchamiani w Q4\_K\_M. W przeciwieństwie do zadań takich jak generacja tekstu, gdzie wpływ kwantyzacji na jakość tekstu czy jego stronniczość jest zjawiskiem zbadanym \cite{marchisioetal2024a,ferrara2026}, tak wpływ kwantyzacji na działanie otwartego modelu jako sędziego pozostaje niszą badawczą \cite{guetal2026a,angheletal2026}.

Znaczna większość badań nad LLM-as-a-judge pozostaje przeprowadzana w oparciu o dane anglojęzyczne \cite{dogruozetal2026a}. Badania tego podejścia w innych językach, zarówno niskozasobowych (ang. \emph{low-resource languages}) \cite{dogruozetal2026a,zhouetal2026}, jak i w języku polskim \cite{karpetal2026}, pozostają natomiast stosunkowo nową dziedziną.

\section*{Wnioski płynące z przeglądu literatury}

Wrażliwość sędziów LLM na kolejność prezentowanych odpowiedzi (ang. \emph{position bias}) jest jednym z najlepiej udokumentowanych zjawisk w tej dziedzinie \cite{zhengetal2023c,wangetal2024,shietal2025a,kooetal2024a}. Zjawisko to nie jest przypadkowe, a jego siła zależy od sędziego i zadania. Najsilniej widać je przy porównywaniu odpowiedzi o zbliżonej jakości \cite{shietal2025a}. Niestety nie da się go wyeliminować prostymi metodami, takimi jak zwiększenie rozmiaru modelu \cite{kooetal2024a}. Autorzy badań, chcąc temu zjawisku zaradzić, proponują zarówno metody jego ograniczania, np. ocenę w obu kolejnościach \cite{wangetal2024}, jak i metryki pomiaru: PC/PF/RS \cite{shietal2025a} oraz RR/CR w ramach frameworku CALM \cite{ICLR2025_fdca08d3}. W kontekście analizowanych w tej pracy problemów szczególnie przydatna wydaje się metryka RS, pozwalająca odróżnić bias od szumu, co przy porównaniu sędziów skwantyzowanych może umożliwić rozdzielenie tych dwóch efektów. Omówione badania przeprowadzono w dużej mierze na modelach zamkniętych i wyłącznie na danych anglojęzycznych, a żadne z nich nie uwzględniało kwantyzacji sędziego, co stanowi lukę badawczą, którą niniejsza praca ma częściowo wypełnić.

Position bias nie jest jednak jedynym problemem sędziów LLM. Jednym z innych biasów jest preferowanie dłuższych odpowiedzi \cite{zhengetal2023c,parketal2024a}, którego działanie okazuje się silnie zależne od modelu, a kierunek bywa przeciwny w różnych rodzinach modeli \cite{soumik2026}. W związku z tym trudno znaleźć na niego jedno uniwersalne rozwiązanie. Innym, często silniejszym efektem jest bias stylu, przejawiający się m.in. preferowaniem odpowiedzi sformatowanych w markdownie \cite{soumik2026}. Osobną grupę stanowi faworyzowanie własnych odpowiedzi (ang. \emph{self-preference}), powiązane ze zdolnością modelu do rozpoznawania własnych tekstów \cite{panicksseryetal2024b} oraz z preferowaniem tekstów, które są dla modelu bardziej przewidywalne, niezależnie od ich autorstwa \cite{wataokaetal2025}. Warto dodać, że silniejsze modele nie wykazują mniejszego self-preference \cite{yangetal2026a}. Gdy natomiast ocenia się treść pozbawioną cech stylistycznych, efekt ten zanika, a oceny zmienia już sama etykieta autorstwa \cite{chaeetal2026a}. W związku z tym w badaniu odpowiedzi powinny być anonimizowane, a pomiar biasu nie powinien ograniczać się do długości tekstu.

Istotnym problemem jest także rzetelność samego pomiaru. Przy identycznych wywołaniach werdykt sędziego (w badaniu dwóch modeli GPT) zmienia się średnio w 13,6\% przypadków, a odtworzenie stabilnego werdyktu wymaga kilkunastu powtórzeń. Zjawisko to częściowo wynika z niedeterministycznego charakteru generacji, jednak dekodowanie deterministyczne jedynie je ogranicza, nie eliminując go całkowicie \cite{yagubyan2026}. Jest to poważny problem, mogący podważyć sens stosowania takich rozwiązań w zastosowaniach krytycznych, w których nie można dopuścić do różnych ocen tego samego przypadku. Powszechnie stosowana surowa zgodność werdyktów sędziego z ocenami ludzi zawyża jego jakość w porównaniu ze współczynnikiem $\kappa$ Cohena, uwzględniającym zgodność przypadkową, a wysoka powtarzalność ocen może współwystępować z silnym biasem \cite{normanetal2026}. Zmiana sędziego, nawet na większy model tej samej rodziny, zmienia z kolei wynik pomiaru \cite{yangetal2026}. Na tej podstawie można stwierdzić, że powtarzanie pomiarów, raportowanie wartości $\kappa$ oraz stosowanie korekty statystycznej wyników są niejako koniecznością \cite{lee2026how}.

Wpływ kwantyzacji na jakość modeli zbadano dotychczas głównie w kontekście generacji tekstu. Jej negatywne skutki są tam znacznie silniejsze, niż wskazują metryki automatyczne: dla języka japońskiego spadek o 1,7\% w testach automatycznych odpowiadał spadkowi o 16\% w ocenie ludzi, przy czym bardziej podatne okazały się mniejsze modele \cite{marchisioetal2024a}. Kwantyzacja może również zwiększać stronniczość odpowiedzi otwartych, której nie wykrywają standardowe testy bezpieczeństwa, przy czym to, czy rośnie ona wraz ze stopniem kompresji, zależy od sędziego oceniającego odpowiedzi \cite{ferrara2026}. Wykorzystanie skwantyzowanych modeli w roli sędziów jest natomiast zagadnieniem nowym. Jedyna znaleziona praca tego typu stosowała jeden poziom kwantyzacji (Q4\_K\_M), co uniemożliwia oddzielenie jej wpływu od cech samego modelu, a proponowane w naszej pracy porównanie poziomów kwantyzacji autorzy wskazują jako kierunek dalszych badań \cite{angheletal2026}.

Badania nad LLM-as-a-judge są zdominowane przez język angielski. Spośród 650 prac z ACL Anthology tylko 33 dotyczyły konfiguracji wielojęzycznych lub języków niskozasobowych (ang. \emph{low-resource languages}), przy czym zwykle wykorzystywano jednego sędziego i rzadko weryfikowano jego oceny przez ludzi \cite{dogruozetal2026a}. Sędziowie osiągają również różne wyniki w zależności od języka, np. gorsze dla języków afrykańskich niż europejskich, a przy porównaniu odpowiedzi w różnych językach preferują odpowiedzi w języku angielskim. Język polski nie był jednak uwzględniony w tych badaniach \cite{zhouetal2026}. Nieliczne prace dotyczące języka polskiego wskazują natomiast, że oceny sędziów LLM mogą znacząco odbiegać od ocen ekspertów \cite{karpetal2026}, choć dużą rolę może tu odgrywać specjalistyczny, mało typowy dla modeli ogólnego przeznaczenia charakter zadań.

W przeanalizowanej literaturze nie znaleziono badań nad wpływem poziomu kwantyzacji otwartego modelu pełniącego rolę sędziego na jego biasy, w szczególności dla języka polskiego, stąd decyzja o uwzględnieniu tego zagadnienia w pytaniach badawczych.

\section*{Pytania badawcze i wstępne przypuszczenia}
Na podstawie dotychczasowych studiów literaturowych oraz wiedzy autorów formułujemy trzy pytania badawcze. Każde z nich dotyczy innej właściwości: P1 cech samego modelu, P2 sposobu jego uruchomienia, a P3 ocenianych danych.

\begin{enumerate}[label=\textbf{P\arabic*.}, leftmargin=2.5em]
  \item Jak siła position bias sędziego LLM zależy od rodziny i rozmiaru modelu?
  \item W jakim stopniu kwantyzacja wag (BF16 $\rightarrow$ Q8 $\rightarrow$ Q4) wpływa na trafność oraz stronniczość (position bias, preferencję dłuższych odpowiedzi) modelu pełniącego rolę sędziego?
  \item Jak język ocenianych odpowiedzi (polski vs angielski) wpływa na trafność i spójność werdyktów sędziów LLM?
\end{enumerate}

Główną motywacją stojącą za sformułowaniem powyższych pytań badawczych są wskazane w trakcie przeglądu literaturowego luki badawcze, dotyczące między innymi wpływu kwantyzacji na działanie modelu jako sędziego oraz oceny w języku polskim.

Na podstawie dotychczasowej wiedzy autorów i wykonanego przeglądu literatury formułujemy następujące przypuszczenia:
\begin{enumerate}[label=\textbf{H\arabic*.}, leftmargin=2.5em]
  \item W obrębie jednej rodziny mniejsze modele wykazują silniejszy position bias niż większe.
  \item Wraz ze spadkiem precyzji wag spada trafność sędziego i rośnie jego stronniczość, przy czym efekt staje się wyraźny dopiero poniżej poziomu Q8.
  \item Sędziowie są mniej trafni i mniej spójni przy ocenie odpowiedzi w języku polskim niż w angielskim, a różnica jest mniejsza dla modelu specjalizowanego w języku polskim.
\end{enumerate}

\section*{Planowane zadania związane z eksperymentem}
W ramach eksperymentów planowane jest wykorzystanie modeli open-weight uruchamianych lokalnie. Względem wstępnych założeń zakres rozmiarów rozszerzono w dół do ok.~4B, co pozwala porównać dwa rozmiary w obrębie każdej rodziny. Wykorzystane zostaną następujące rodziny modeli (w nawiasach rozmiary): Qwen3.5 (4B, 9B), Gemma~4 (E4B, 12B), Ministral~3 (8B, 14B) oraz Bielik~v3 (4.5B, 11B), jako model specjalizowany w języku polskim. Opcjonalnie, jako punkt odniesienia często spotykany w literaturze, uwzględniony zostanie Llama~3.1 (8B). We wszystkich modelach tryb rozumowania (ang. \emph{thinking}) pozostanie wyłączony.

Eksperymenty zostaną przeprowadzone następująco:
\begin{itemize}[leftmargin=1.5em]
  \item \textbf{P1:} wszystkie wymienione modele w kwantyzacji Q4\_K\_M, tak aby kwantyzacja pozostała stała. Każda para odpowiedzi oceniana jest dwukrotnie, w kolejności A--B oraz B--A. Mierzone będą spójność werdyktu po zamianie kolejności, odsetek wyborów pierwszej pozycji oraz zgodność z ocenami ludzi (współczynnik $\kappa$ Cohena).
  \item \textbf{P2:} dwa modele z różnych rodzin (Qwen3.5-9B oraz Ministral-3-8B) uruchamiane na kolejnych poziomach precyzji: BF16 (wagi niekwantyzowane), Q8\_0, Q6\_K, Q4\_K\_M i Q3\_K\_M. Wszystkie poziomy zostaną utworzone samodzielnie z tego samego pliku BF16 przy użyciu jednego narzędzia, a w raporcie podana zostanie efektywna liczba bitów na wagę. Mierzone będą trafność na parach z obiektywnie poprawną odpowiedzią oraz zmiana position bias i preferencji dłuższych odpowiedzi między poziomami.
  \item \textbf{P3:} wybrane modele z P1 oceniają te same pary odpowiedzi w wersji angielskiej i polskiej. Porównana zostanie trafność oraz spójność werdyktów między językami.
\end{itemize}

Wszystkie modele zostaną uruchomione przy użyciu jednego silnika inferencji (llama.cpp) w jednej, ustalonej wersji, w formacie GGUF. Dzięki temu porównywane poziomy kwantyzacji nie różnią się dodatkowo implementacją. Odpowiedź sędziego zostanie ograniczona gramatyką do jednego z dwóch werdyktów, a poza samym werdyktem zapisywane będzie prawdopodobieństwo wyboru każdej z odpowiedzi, co pozwoli wykryć również niewielkie zmiany w zachowaniu sędziego. Główny pomiar przeprowadzony zostanie przy dekodowaniu deterministycznym, a dodatkowe powtórzenia przy niezerowej temperaturze posłużą do oszacowania stabilności ocen.

Eksperymenty zostaną przeprowadzone na trzech maszynach: dwóch komputerach z kartą RTX~5080 (16~GB VRAM, 64~GB RAM) oraz komputerze MacBook Pro z procesorem M1 (16~GB pamięci współdzielonej). Każdy eksperyment wykonywany będzie w całości na jednej maszynie, aby różnice w obliczeniach między platformami nie wpływały na porównania. Modele w precyzji BF16, które nie mieszczą się w całości w pamięci karty graficznej, zostaną częściowo przeniesione do pamięci RAM, co wydłuża czas obliczeń, ale nie zmienia ich wyniku.

Opcjonalnie, jako punkt odniesienia, planowane jest wykorzystanie modelu zamkniętego dostępnego przez API, przy czym ograniczeniem będzie jego koszt. Jeżeli model taki posiada tryb rozumowania, w trakcie eksperymentu pozostanie on wyłączony.

\section*{Dane}
\begin{itemize}[leftmargin=1.5em]
  \item \textbf{MT-Bench human judgments} \cite{zhengetal2023c}: pary odpowiedzi z ocenami ludzi, standardowy zbiór do badania position bias (P1).\\
  \url{https://huggingface.co/datasets/lmsys/mt_bench_human_judgments}
  \item \textbf{LLMBar} \cite{ICLR2024_afc8b034}: 419 trudnych par, w których gorsza odpowiedź wygląda lepiej (P1, P2, P3).\\
  \url{https://github.com/princeton-nlp/LLMBar}
  \item \textbf{JudgeBench} \cite{ICLR2025_9e720fce}: pary z obiektywnie poprawną odpowiedzią (wiedza, rozumowanie, matematyka, kod) (P2).\\
  \url{https://github.com/ScalerLab/JudgeBench}
  \item \textbf{EvalBiasBench (OffsetBias)} \cite{parketal2024a}: 80 ręcznie przygotowanych przypadków na konkretne typy biasu (P2).\\
  \url{https://github.com/ncsoft/offsetbias}
\end{itemize}

\subsection*{Dane w języku polskim}
\begin{itemize}[leftmargin=1.5em]
  \item \textbf{LLMBar-PL}: polskie tłumaczenie zbioru LLMBar, przygotowane w ramach projektu. Etykieta wskazująca odpowiedź zgodną z poleceniem zachowuje ważność po tłumaczeniu, dzięki czemu możliwe jest porównanie trafności sędziów na tych samych parach w obu językach (P3). Jakość tłumaczenia zostanie zweryfikowana ręcznie na losowej próbie.
  \item \textbf{PLLuM-Align}: polski zbiór preferencji tworzony w całości przez ludzkich anotatorów (oceny odpowiedzi, rankingi, dialogi wieloturowe), wykorzystany jako dodatkowy zbiór natywnie polski (P3).\\
  \url{https://huggingface.co/datasets/NASK-PIB/PLLuM-Align}
  \item \textbf{MT-Bench Polish (lightblue)}: 80 pytań MT-Bench przetłumaczonych na polski (GPT-4o, ręczna korekta przez rodzimą użytkowniczkę). Zbiór nie zawiera odpowiedzi ani ocen, dlatego może posłużyć jedynie do wygenerowania nowych par odpowiedzi.\\
  \url{https://huggingface.co/datasets/lightblue/mt_bench_polish}
\end{itemize}


\printbibliography[title=Bibliografia]

\end{document}